\scanpage{256}

\[1\longrightarrow\underset{\substack{\shortparallel\\Z}}{T_H}\longrightarrow T_G\xrightarrow{\widetilde\delta|T_G}\widetilde Z\longrightarrow1\tag*{(117)}\]

qui est scindée par $T_{0G}$ et $T_{1G}$, qui donnent deux \deleted{\unsure{iso}} hom

\diagram{\begin{tikzcd}[ampersand replacement=\&]
(118)\quad\widetilde Z \arrow[r,shift left=2.6pt,"{\deleted{\unsure{\delta|T_0}}u_0}"] \arrow[r,shift right=2.6pt,swap,"u_1"] \& T_G
\end{tikzcd}}

satisfaisant la relation

\[\begin{aligned}
\underset{\substack{s_0\mapsto\deleted{\unsure{\widetilde{\phantom{\lambda}}}}\lambda s_0\\t_0\mapsto t_0}}{u_1(\widetilde\lambda)}\deleted{\unsure{u_0}}
={}&\underbrace{u_0(\widetilde\lambda)}_{\substack{(s_0\mapsto s_0)\\(t_0\unsure{-|\mapsto}\lambda t_0)}}
\deleted{\unsure{u_{c(\widetilde\lambda)}}}
\underbrace{\overbrace{u_{c(\widetilde\lambda)}}^{\text{él. central de }T_G}}_{\substack{s_0\mapsto\lambda s_0\\t_0\mapsto\lambda^{-1}t_0}}
\end{aligned}\tag*{(119)}\]

qui \deleted{\unsure{se}} s'écrit aussi

\[\begin{aligned}
u_1(\widetilde\lambda)u_0(\widetilde\lambda)^{-1}&=u_{c(\widetilde\lambda)}\\
\text{ou }\ u_0(\widetilde\lambda)^{-1}u_1(\widetilde\lambda)&
\end{aligned}\]

et comme les \deleted{\unsure{premiers}} deux mbes sont dans $H$, il suffit de vérifier que les valeurs dans $s_0,t_0$ sont les $=$s, ce qui est immédiat.

L'un et l'autre scindage de (117) donnent deux isomorphismes

\diagram{\begin{tikzcd}[ampersand replacement=\&]
(120)\quad {\deleted{\unsure{Z}}T_G} \arrow[r,shift left=2.5pt,"\sim"] \arrow[r,shift right=2.5pt,"\sim" description] \& Z\times\widetilde Z,
\end{tikzcd}}

le passage de l'un à l'autre se faisant via l'hom $\widetilde Z\xrightarrow{k}Z$.

Quand $\widetilde Z$ est commutatif, \inserted{donc $T_G$ commutatif,} cette description se simplifie, en écrivant la suite exacte

\[\begin{aligned}
0\longrightarrow\unsure{\mathfrak{z}|\mathfrak{Z}}\longrightarrow
&\overbrace{T_0\unsure{\mathbin{.}\times|\times}T_1}^{\simeq\widetilde Z\times\widetilde Z}\longrightarrow T_G\longrightarrow1\\
&(t,t')\longmapsto t.t'\\[-2pt]
z\longmapsto (z,z^{-1})&
\end{aligned}\tag*{(121)}\]

qui explicite $T_G$ comme un quotient de $T_0\times T_1\simeq\widetilde Z\times\widetilde Z$

% archive sheet number: 255
